\section{The holomorphic theta family}\label{sec:theta}

We construct forms on a unitary group of signature $(3,1)$ at each real
place of $F$, with unit eigenvalues for two normalized positive
parabolic operators at $p$ and coefficients of uniformly bounded
denominator.  These are the forms
whose constant and positive Fourier--Jacobi coefficients will be compared
in the next two sections.  The operator calculation allows arbitrary
reduction of $A_0$ at $p$; the denominator bound must hold as both the
weight and the tame level grow.  Throughout this section a finite-order
character at the moving prime means the universal character with values
in a group algebra.  All
integrality assertions are made before specializing that algebra at its
characters.

\subsection{Hermitian spaces and the compact automorphic input}

Let $\delta=\sqrt D$, with $\operatorname{Im}\sigma(\delta)>0$ for
$\sigma\in\Sigma$.  We use the Hermitian space and the skew-Hermitian space
\begin{equation}\label{eq:theta-spaces}
 A=(L^2,1_2),\qquad
 W=B\oplus\mathbb H,\qquad
 B=(L^2,b_0\delta\,1_2),
\end{equation}
where $\mathbb H=Le\oplus Lf$ is a hyperbolic plane and
$b_0\in F^{\times,+}$ is a unit at the primes above $p$.  Choose $b_0$ so that
\begin{equation}\label{eq:lattice-parities}
 v(b_0\delta\mathfrak d_F)\equiv0\pmod2
 \quad\text{at every nonsplit finite place of }L/F.
\end{equation}
Here $\mathfrak d_F$ is the different and we use the trace additive
character.  To check the existence of this choice, prescribe the local
norm classes of $b_0$ needed to cancel the parity of the additive conductor
on each paired line.  The exact norm-invariant sequence for the quadratic
extension $L/F$ says that these prescriptions come from a global element
if and only if their product is one.  The possible odd contributions are
at the primes of $D'$ that are inert in the relevant quadratic extension;
the condition $(D/|D'|)=1$ in Lemma~\ref{lem:auxiliary-fields} makes their
number even.  The discriminants are odd and coprime, and the places over
$2$ are split.  The real prescriptions are positive.  The product condition
therefore holds.  Multiplication by a global norm, using approximation at
the split $p$-places, makes $b_0$ a $p$-unit without changing the prescribed
classes.

At a nonsplit finite place the extension is unramified.  Scaling a line
coefficient by a norm changes its valuation by an even integer, so
\eqref{eq:lattice-parities} gives a self-dual Schr\"odinger lattice for each
paired line, compatible with the trace character.  Choose compatible
hyperbolic bases.  These choices give the unramified models and spherical
tests used below.  At $p$ all coordinates are unit self-dual coordinates.
We pair the sign of the alternating form with the additive character so
that a positive unary theta series has exponential
$\exp(2\pi i a x^2z)$ with $a>0$.

Write $\pi$ for the representation on the compact group $U(A)$ obtained
from the weight-two form of $A_0$ by real quadratic base change and
Jacquet--Langlands transfer, followed by pullback to $U(A)$; see
\cite{ArthurClozel1989,JacquetLanglands1970}.  More explicitly,
$PU(A)$ is the projective multiplicative group of the quaternion algebra
ramified at the two real places of $F$ and at no finite place.  The real
component of $\pi$ is trivial, as is its adelic central character.
The successive base changes remain cuspidal: the only possible
quadratic inducing field in the dihedral case was excluded in the choice
of $L$.  Its standard base change to $L$, denoted $\pi_L$, is thus the
cuspidal base change of the representation of $A_0$.

At a nonsplit finite place take hyperspecials with the $B$-lattice
homothetic to the $A$-lattice.  The hyperspecial maps onto the hyperspecial
of the adjoint quotient because the unramified kernel is connected.  On
dual groups the quotient map is the inclusion
$\mathrm{SL}_2\hookrightarrow\mathrm{GL}_2$ with trivial pinned Galois
action on the former.  Consequently the spherical standard parameters
are the restrictions of those of the quaternionic transfer.  We use
algebraic-valued vectors in the compact automorphic model of $\pi$;
a rational tensor vector can be taken real algebraic.  They are spherical
at nonsplit finite places.  No local condition is imposed on $\pi_p$.
The precise multiplicity and pullback assertion needed for the binary
lift is proved in Section~\ref{sec:constant-terms}.

For the moving prime $q_i$, let $\nu_i$ denote the universal character of
$\Delta_i$, pulled back to the norm-one torus of $L/F$ via the norm
$N_{L/K}$.
Put
\begin{equation}\label{eq:theta-character}
 \eta=\cC^m\nu_i,\qquad k_0=m-1.
\end{equation}
The integer $m>0$ will grow in the ordinary sense and tend to zero
$p$-adically on the fixed torsion branch.  The finite smooth conductor
of the factor $\cC^m$ is fixed, or is killed by that branch.
The tame character $\nu_i$ is trivial
at every fixed place where this is required by
Proposition~\ref{prop:tame-primes}.  We often suppress $i$ in $\eta$.

\subsection{The Fock vector and its algebraic weight}

Take the oscillator representation on $A\otimes_L W$ with both even
unitary splitting characters trivial, in the conventions of
\cite{Kudla1994}.  For a bare compact vector $z\in\pi^\vee$, the theta
integral has the form
\begin{equation}\label{eq:upper-theta-integral}
 \Theta_{i,m}(g;\Phi,z)
 =\int_{[U(A)]}\theta_\Phi(g,h)z(h)\eta^{-1}(\det h)\,dh.
\end{equation}
The quotient of integration is compact, so this is an ordinary
convergent theta integral.

At each real place use Fock polynomials of degree $k_0$ in the
$2\times2$ minors of the $2\times3$ positive-coordinate matrix.  Their
space is
$\operatorname{Sym}^{k_0}(\bigwedge^2\operatorname{Std}_3)$.
The vacuum contributes $\det$ on the source $U(2)$, while each minor
contributes another determinant.  Thus the source type is $\det^m$, as
required by \eqref{eq:upper-theta-integral}.  On the target the vacuum has
weight $(1,1,1;-1)$, and the polynomial contributes
$(m-1,m-1,0;0)$.  The resulting weight is
\begin{equation}\label{eq:theta-weight}
 \lambda=(m,m,1;-1).
\end{equation}
The first three entries refer to the distinguished holomorphic Lie
quotient; the last refers to the remaining distinguished direction,
whose conjugate is holomorphic.  The antiholomorphic lowering operators
are contractions involving the negative coordinates, which do not occur
in these polynomials.  They therefore annihilate the vector.  The
mixed expansion below proves holomorphy also at the cusps.

We specify the integral coefficient lattice, including its frame
convention.  Let $\mathcal H=\omega^+$ be the rank-three positive
Hodge differential bundle and put $d=m-1$.  The positive creation
directions belong to the holomorphic Lie quotient, dual to
$\mathcal H$.  Write frames as row tuples.  If a differential frame
changes from $e$ to $eg$,
its dual creation frame changes by $g^{-t}$, and the minor frame changes
by $\bigwedge^2g^{-t}$.  This passive change of differential frame is
distinct from the active compact-group action on the Fock polynomials
used to compute \eqref{eq:theta-weight}: the Lie quotient has active
type $\operatorname{Std}_3$, while its differential dual has type
$\operatorname{Std}_3^\vee$.

Theta evaluation is linear in its Fock test polynomial.  Equivalently,
the covariance identity
$\Theta(gk;P)=\Theta(g;\omega(k)P)$ says that its coefficient array
transforms dually to the input polynomial lattice.  In the preceding
differential frames this gives
\begin{equation}\label{eq:dual-fock-lattice}
 \mathcal P_d(\mathcal H)
   =\operatorname{Sym}^d((\textstyle\bigwedge^2\mathcal H)^\vee),
 \qquad
 \mathcal D_d(\mathcal H)
   =\mathcal P_d(\mathcal H)^\vee
   =\Gamma^d(\textstyle\bigwedge^2\mathcal H).
\end{equation}
Here $\Gamma^d$ denotes the divided-power module, defined by the
displayed duality for finite locally free modules.  We use the literal
dual basis to the Fock monomials throughout.  In particular, extracting
a displayed monomial coefficient does not multiply it by a binomial
coefficient or a factorial.  Including the vacuum, the positive
coefficient bundle is
$\det(\mathcal H)\otimes\mathcal D_d(\mathcal H)$, and the
opposite differential line occurs to exponent $1$.
Section~\ref{sec:cusp-lifting} uses this same integral bundle.
All constructions are tensored over the two embeddings of $F$.

We use the PEL model with common rational similitude to give these
coefficient bundles their algebraic meaning, and restrict to the
isometry components.  Analytically these are the Hermitian symmetric
domain times finite isometry-group cosets modulo the rational isometry
group.  Fixing the multiplier class of the finite frame and absorbing
positive principal multipliers into the polarization realizes them as
components or finite covers of the PEL model, with compatible cyclotomic
constants.  The positive rational units preserving a fixed compact
multiplier level are trivial.  Thus the isometry Hecke correspondences
and the expansion principles below are compatible with this model.
No extension of the theta kernel to a chosen similitude theta
representation is required.

\subsection{Split finite tests and the two parabolic operators}

At a split finite place identify the isometry groups with their actions
on the distinguished factors.  In the linear-pair polarization the
Schr\"odinger variables are dual covectors $Y\in M_{2,4}$, and the action
is
\begin{equation}\label{eq:split-weil-action}
 \omega(h,g)\Phi(Y)
  =|\det h|^2|\det g|\Phi({}^thYg).
\end{equation}
For the binary target the exponent on $|\det h|$ is $1$ instead of $2$.
There is no further determinant character: both even splitting
characters are trivial.  At nonsplit unramified places the matched
self-dual lattice test is spherical.  At every other fixed finite place
we use integral tests up to a fixed denominator and at fixed level.
Finite additive Fourier transforms away from $p$ introduce only powers
of the residue characteristic, roots of unity, and square roots of
moduli.  These are $p$-units; in particular no inverse order of a
multiplicative residue group is introduced.

Order the distinguished columns at $p$ as $(e,B_1,B_2,f)$.  The $e$
direction is multiplicative.  Our tests have the form
\begin{equation}\label{eq:p-pivot-test}
 \Phi_p(Y)=\phi_p(Y_{[1,2]})
           \1_{\Z_p^2}(Y_3)\1_{\Z_p^2}(Y_4),
 \qquad \supp\phi_p\subset\GL_2(\Z_p),
\end{equation}
where $\phi_p$ is locally constant at a fixed depth $s$.
They are fixed by a subgroup $J\subset\GL_4(\Z_p)$ with upper block
pattern $(2,1,1)$: its upper block unipotent is full, its first Levi is a
principal congruence subgroup of depth $s$, its two remaining unit factors
are full, and its lower block entries have depth $s$.  Increasing the
fixed $s$ if needed makes the same description work for all the tests.
On the canonical component the rank-three distinguished flag member
modulo $p^s$ is the multiplicative subgroup.  This level records only
flags and frames in that subgroup; it does not choose a splitting of the
\'etale quotient.  Its integral geometric description will be given in
Section~\ref{sec:cusp-lifting}.

At either distinguished $p$-place put
\begin{equation}\label{eq:theta-u-operators}
 t_j=\diag(1_j,p^{-1}1_{4-j}),\qquad U_j=[Jt_jJ],\qquad j=2,3,
\end{equation}
where $U_j$ is the unnormalized right-coset sum.

\begin{lemma}\label{lem:split-u-operators}
 Every test \eqref{eq:p-pivot-test} satisfies $U_2\Phi_p=p^2\Phi_p$
and $U_3\Phi_p=p^2\Phi_p$.  Thus both operators $p^{-2}U_j$ have
value $1$ on the theta form of weight \eqref{eq:theta-weight}.
\end{lemma}
\begin{proof}
 The upper shifts in the double coset are indexed by
 $M_{j,4-j}(\F_p)$, of cardinality $p^{j(4-j)}$.  The expanding columns
 must become divisible by $p$ before the multiplication by $p^{-1}$.
 Since the first two columns have rank two modulo $p$, these equations
 have $p^{(j-2)(4-j)}$ solutions whenever the original test is supported,
 and none otherwise.  The equations do not change the chosen head
 function.  The determinant factor in \eqref{eq:split-weil-action} is
 $p^{4-j}$.  Hence the surviving factor is
 \[
 p^{4-j}p^{(j-2)(4-j)}=p^{(j-1)(4-j)}=p^2
 \]
 for both $j=2$ and $j=3$.  The calculation is independent of the source
 representation and of its reduction at $p$.
\end{proof}

At each of the two places of $F$ over $q=q_i$ take the corresponding
pivot test
\begin{equation}\label{eq:q-pivot-test}
 \Phi_q(Y)=
 \1_{\GL_2(\cO_q)}(Y_{[1,2]})
 \eta_q(\det Y_{[1,2]})
 \1_{M_{2,2}(\cO_q)}(Y_{[3,4]}).
\end{equation}
Only the tame determinant factor is relevant on the displayed compact
support.  Its transformation under $\GL_2(\cO_q)$ cancels the
$\eta_q^{-1}$ in source integration.  Consequently this integration uses
a compact subgroup of fixed volume, rather than the kernel of a
character of growing order.  The coefficients in
\eqref{eq:q-pivot-test} are group-like units in the tame group algebra.
The corresponding two-column test is the reference test in the binary
kernel.  Section~\ref{sec:content} will impose the other finite local
controls without changing either \eqref{eq:p-pivot-test} or
\eqref{eq:q-pivot-test}.

\begin{lemma}\label{lem:theta-eigenparameters}
 At a split spherical place the standard parameter of the theta form,
 in unitary normalization, is
 \begin{equation}\label{eq:theta-eigenparameters}
 (\pi\otimes\eta)\ \oplus\ |\cdot|^{-1/2}
                       \ \oplus\ |\cdot|^{1/2}.
 \end{equation}
 Here $\pi\otimes\eta$ means tensoring the standard two-dimensional
 parameter by the determinant-twist character.
\end{lemma}
\begin{proof}
 This is the unramified theta correspondence for the matched lattices.
 At split places it can also be read directly from the open rank-two
 matrix orbit: its coinvariants are the normalized parabolic induction
 of the determinant-twisted source with the trivial binary
 representation.  The latter has standard parameters
 $|\cdot|^{-1/2},|\cdot|^{1/2}$.  The open-orbit calculation, including
 the exclusion of the lower ranks near the comparison point, is given
 in Lemma~\ref{lem:fixed-place-spanning}.
\end{proof}

The spaces, coefficient lattice, and local tests are now fixed.  We
record the resulting family before proving the uniform denominator
bound, which occupies the rest of this section.

\begin{proposition}\label{prop:theta-family}
 For the characters \eqref{eq:theta-character} and every sufficiently
 large positive $m$ on the prescribed branch, the theta construction
 gives holomorphic algebraic forms $\Theta_{i,m}$ of weight
 \eqref{eq:theta-weight}.  In the canonical ordinary monomial-dual
 lattice their coefficients have a common bounded denominator,
 coefficientwise in the tame group algebra and in every fixed
 truncated substitution $R_b$.  They have spherical parameters
 \eqref{eq:theta-eigenparameters}, and at both distinguished $p$-places
 they satisfy
 \[
  p^{-2}U_2\Theta_{i,m}=\Theta_{i,m},\qquad
  p^{-2}U_3\Theta_{i,m}=\Theta_{i,m}.
 \]
 These assertions allow arbitrary fixed split tests of the stipulated
 bounded levels.  They remain valid for the particular tests with all
 the finite local controls in Proposition~\ref{prop:bounded-content}.
\end{proposition}

\subsection{Frames, periods, and unary moments}

The remaining bound reduces, through the mixed Fourier--Jacobi
expansion, to moments of unary theta series.  We first fix their periods
and differential frames.  The essential estimate will bound every
polynomial combination of moments by the supremum norm of that
polynomial on the position lattice; bounds for individual degrees
alone would not give the required interpolation.

Use the PEL Hodge bundles, with logarithmic differentials in toric
directions.  In complex linear coordinates analytic-to-de Rham comparison
uses $2\pi i\,du$ for each holomorphic differential, including those of
the conjugate Lie quotient.  Normalize each Fock creation generator so
that its leading Schr\"odinger position coordinate is linear, rather
than multiplied by $\pi^{1/2}$.  Diagonal CM and hyperbolic frames can be
chosen with algebraic $p$-unit scales, allowing square roots.  The base
heights, quadratic coefficients, and comparison bases can likewise be
chosen to be units at $p$.

On one positive paired line, restriction to the symplectic
$\mathrm{SL}_2/F$ gives a Hilbert unary theta function at the CM point
of type $\Sigma$.  Write $d=e_0+2j$ with $e_0\in\{0,1\}$, independently
at each embedding when necessary.  The $d$-th Fock vector is the
$j$-fold normalized Maass--Shimura raise of the base theta of weight
$\frac12+e_0$.  Indeed the compact-circle raising operator multiplies by
the square of the Fock variable, while the leading term of the normalized
derivative of the Gaussian exponential is its quadratic form.
Comparison of leading terms identifies the two vectors up to a fixed
algebraic unit raised to the degree.  There is no factor $j!$.

If $\Omega_\sigma$ compares the linear differential with the chosen
integral algebraic CM differential, normalization multiplies by
\begin{equation}\label{eq:unary-period-normalization}
 \left(\frac{2\pi i}{\Omega_\sigma}\right)^{1/2+d}.
\end{equation}
Use compatible square roots in products.  Since the type is induced from
$K$, rational tensor decompositions of homology over $K$, with bases
localized at $p$, compare these periods with the elliptic CM period by
algebraic $p$-units.  Equivalently use the Serre tensor construction and
prime-to-$p$ isogenies; $p$ is completely split.  Passing to canonical
Igusa differentials gives the analogous comparison with $p$-adic CM
periods, again with unit factors.

\begin{convention}\label{conv:scaled-support}
 The position-moment functional is defined first in fixed Igusa
 differential frames.  Fixed powers of algebraic unit scales arising
 from conversion to another CM frame are restored afterwards.  An
 equivalent formulation puts the polynomial norm on the actual scaled
 compact support.  We do not identify that support with $\Z_p$ merely
 because the scale is an algebraic $p$-unit: a unit in a coefficient
 extension need not belong to $\Z_p$.  The restored unit powers are
 analytic on a sufficiently narrow weight branch and have valuation
 zero.
\end{convention}

\begin{lemma}[Uniform unary moments]\label{lem:unary-moments}
 Fix the integral position supports, the locally constant tests at $p$,
 and their levels.  At a fixed ordinary CM test object, the
 period-normalized unary jets, in the frames of
 Convention~\ref{conv:scaled-support}, are the moments of a bounded
 measure on the product of their integral position lattices.  Its norm
 is bounded in terms of the coefficients of the fixed tests and their
 $p$-levels.  The bound is independent of the moment degrees and of the
 moving tame level $q_i$.

 Cutting a position measure to units gives at classical degrees the
 theta jet with the corresponding unit cut on the finite position test.
 The corresponding moments converge to these unit moments when every
 degree tends to infinity and converges on its fixed $p$-adic branch.
 On a
 sufficiently narrow branch the unit moments, including the fixed CM
 unit-scale powers, are analytic functions of the degrees.
\end{lemma}
\begin{proof}
 For a scalar position test $\phi$, write $J_{\mathbf d}$ for the
 normalized jet of multi-degree $\mathbf d$ at the fixed CM object.
 For a polynomial
 $P(X)=\sum_{\mathbf d}a_{\mathbf d}X^{\mathbf d}$ on the product
 position lattice $\Lambda$, put
 \[
  L_\phi(P)=\sum_{\mathbf d}a_{\mathbf d}J_{\mathbf d},
  \qquad \|P\|_\Lambda=\sup_{x\in\Lambda}|P(x)|_p.
 \]
 We will prove $|L_\phi(P)|_p\leq C\|P\|_\Lambda$, where $C$
 depends only on the stated test bounds and $p$-levels.  The proof
 constructs the square of this combination as an algebraic Igusa
 function, bounds its expansions, and removes its initial denominator.

 The square argument in this proof is scalar.  For a test with values
 in a tame group algebra, first expand it in the group basis and apply
 the argument separately to each scalar coefficient test.  These tests
 have the same coefficient bound and the same position supports and
 levels at $p$.  Only after obtaining the scalar estimates do we
 reassemble the group-algebra-valued moments.  In particular, no
 coefficientwise bound is inferred by taking a square root inside a
 growing group algebra.

 \paragraph{Algebraic squares in a common Igusa ring.}
 For a fixed tame Schwartz level $N_1$, which may contain $q_i$, let
 $\mathcal I_{N_1}$ be the complete ordinary Hilbert Igusa tower with its
 canonical multiplicative trivialization.  In its canonical
 logarithmic differential frames all integral Hodge weights are
 functions.  Let $\mathcal A_{N_1}$ be the $p$-adically complete ring of
 these regular functions.  Finite sums of different weights therefore
 lie in one ring; a finite Igusa level may be chosen separately at each
 reduction precision.  We use Katz's expansion principle
 \citep[Section~5.2]{Katz1976} and its Hilbert formulation
 \citep[Theorem~6.10 and Sections~11.1--11.7]{AndreattaGoren2005}.
 The coefficient-multiplier formula for differentiation is
 \citet[Sections~15.21--15.24]{AndreattaGoren2005}, which identifies
 these operators with those of \citet[Corollary~2.6.25]{Katz1978}.
 For the CM comparison we use the equality of the ordinary unit-root
 and CM Hodge projections in \citet[Section~1.4, Theorem~1.2]{HidaTilouine1993},
 attributed there to Katz~(2.6.7).  The prime $p$ splits
 completely in our CM field, and the chosen multiplicative direction
 defines the required ordinary CM type.  After the fixed tame-level
 enlargement, the level and polarization hypotheses in that comparison
 hold.  We apply it to the algebraic integral-weight products constructed
 below; their algebraization and descent are part of this proof.

 We first show that the square of a finite combination of jets belongs
 to $\mathcal A_{N_1}[1/p]$.  Enlarge the fixed prime-to-$p$
 level, including its level at $2$, to fix the metaplectic multipliers.
 Every product of two base half-weight thetas is then an ordinary
 integral-weight algebraic form, by its Fourier expansion.  Products
 of their normalized nearly holomorphic derivatives are rational
 expressions in the jets of these base products: apply Leibniz and
 Gauss--Manin, initially allowing powers of base squares in the
 denominators.  There are no horizontal poles.  On a smooth
 characteristic-zero Hodge-frame chart the square of a base theta has
 even horizontal divisor, as is seen in its holomorphic half-weight
 realization over $\C$.  Locally write this square as $h^2u$, with
 $u$ invertible.  Adjoining $\sqrt u$ is \'etale, and the half-weight
 coefficient $h\sqrt u$ is regular.  On the torsor of complements to
 the Hodge filtration, projected Gauss--Manin differentiation has a
 regular connection matrix; its half-weight action multiplies the
 scalar weight by $1/2$.  Differentiating $h\sqrt u$ therefore
 introduces powers of $2$ and $u$ in denominators, but no inverse of
 $h$.  All fixed iterates are regular.  Products descend because the
 base cross-products fix the signs of the local square roots.
 In logarithmic cusp charts, choose fixed widths clearing the
 half-integral exponents; the operators $q_\sigma\partial/\partial
 q_\sigma$ preserve nonnegative orders.  The same argument applies
 there, and specializing the splitting variable gives the required
 jets.  Omit an identically zero base theta.

 Projection by the ordinary unit-root splitting consequently produces
 a function $F_P\in\mathcal A_{N_1}[1/p]$ from the square of the
 combination defining $L_\phi(P)$.  Its initial vertical denominator
 may depend on the orders and on $N_1$.  At a CM object the unit-root
 splitting is the CM splitting, so these products agree with the normalized analytic jets,
 including products of different orders and weights.  In particular,
 $F_P(\mathrm{CM})=L_\phi(P)^2$.  Compatible square roots also give
 algebraicity of the individual normalized values.  The common Igusa
 ring is what makes this construction possible when $P$ combines
 different weights.

 \paragraph{Compatible cusps and the position norm.}
 We next choose cusp frames in which the bound by $\|P\|_\Lambda$
 is preserved.  On each paired unary line, choose the position vector
 in the distinguished $L$-factor selected by $\Sigma_p$, the same
 factor used for the linear-pair coordinates $Y$.  This is the CM
 multiplicative eigenspace.  The unit local pairing makes this an
 integral symplectic frame.  Specify the finite position tests in
 this frame from the outset.  At a place over $p$, choose
 $s$ such that all these tests are supported on $\Z_p$ and constant
 modulo $p^s\Z_p$.  Their additive Fourier transforms are supported
 on $p^{-s}\Z_p$.  With the unramified additive character and unit
 quadratic coefficient, the lower unipotent with parameter
 $c\in p^{2s}\Z_p$ fixes each test: conjugating by Fourier transform
 makes its action multiplication by the trivial phase
 $\psi(cQ(x))$ on that support.  Upper unipotents with integral
 parameter preserve integral position support, and diagonal units
 rescale it by units.  Enlarge a fixed depth $M$ to cover all tests
 and fixed character conductors.  Their congruence stabilizer then
 contains the subgroup with $a,d\equiv1$, $c\equiv0\pmod {p^M}$
 and $b\in\Z_p$.

 Use the canonical $\Gamma_0(p^M)$ section over the tame ordinary
 moduli space, together with the multiplicative Igusa marking.
 At a Tate cusp its connected subgroup is $\mu_{p^M}$, the same
 position direction used at the CM object.  Choose cusp
 representatives integral at $p$ with unit diagonal entries, by
 choosing their cusp ideals prime to $p$ and applying approximation
 to the bases.
 Their effective frame changes lie in $B(\Z_p)K(p^M)$, where $B$
 is the upper triangular subgroup and $K(p^M)$ the principal
 congruence subgroup.  The congruence part fixes the finite test;
 the triangular part preserves its support and coefficient norm.
 Thus this choice involves no initial Fourier interchange of the
 position and momentum directions.

 These cusps test every component.  Indeed the finite Igusa cover
 over each connected tame ordinary component is a Galois torsor
 under the units of the multiplicative character group
 \citep[Sections~7.2 and~11.1]{AndreattaGoren2005}.  Its deck group acts
 transitively on the connected components of the cover.  The
 canonical unramified Tate-cusp chart of
 \citet[Sections~6.5 and~11.7]{AndreattaGoren2005} gives one such
 cusp; translating its marking by the deck units gives a cusp on
 each component.  These changes are diagonal units and preserve the
 position norm.  We use these canonical charts and their unit-deck
 translates as the compatible cusps.  The expansion principle applies componentwise, so
 one such cusp per component will suffice to detect vertical
 divisibility.  All of this concerns the position frames before restoring
 the separate unit scales of Convention~\ref{conv:scaled-support}.

 \paragraph{The uniform bound and removal of vertical denominators.}
 By scaling $P$, it suffices to prove the estimate when
 $\|P\|_\Lambda\leq1$.  At a compatible cusp the base expansion is
 $\sum_x\phi(x)q^{Q(x)}$.  The normalized jet of degree $e_0+2j$ has
 expansion
 \[
  \sum_x\phi(x)x^{e_0}Q(x)^j q^{Q(x)}.
 \]
 Katz differentiation multiplies by the quadratic index.  The
 quadratic coefficient is a $p$-unit; normalize the position coordinate
 first, or record its power as a separate fixed unit-scale factor.
 A polynomial $P$ of supremum norm at most $1$ on the integral position
 lattice therefore gives coefficients
 \begin{equation}\label{eq:moment-cusp-coefficient}
  \sum_{Q(x)=\beta}\phi(x)P(x),
 \end{equation}
 each bounded by the fixed bound for $\phi$.  Every sum is finite.
 The square has bounded convolution coefficients, with twice the
 valuation loss.  In several variables the same argument applies to
 multiorders.  Parity pieces are separated by the sign projectors
 $2^{-r}\sum_{\varepsilon\in\{\pm1\}^r}\varepsilon^e[\varepsilon]$,
 which are integral because $p>2$; the base cross-products give their
 common descent data.

 At primes other than $p$, changes of polarization preserve this
 bound, since their normalizations are residue-characteristic powers
 and roots of unity.  At $q_i$ the scalar coefficient tests obtained
 from the determinant character have integral coefficients.  Additive
 Fourier transforms there likewise have $p$-unit normalization.
 Thus $F_P$ has the same expansion bound at every compatible cusp,
 independently of the degree and $q_i$.  Clear that fixed bound and
 call the resulting function $G\in\mathcal A_{N_1}[1/p]$.
 Its expansions are integral.  Choose the least $r\geq0$ for which
 $\varpi^rG\in\mathcal A_{N_1}$.  If $r>0$, all expansions of
 $\varpi^rG$ vanish modulo $\varpi$.  At a finite level containing
 this reduction, the expansion principle on each ordinary component
 implies that its reduction is zero: every component meets one of
 the compatible cusps just constructed, and the finite \'etale
 Igusa cover introduces no denominator.  Thus $\varpi^rG$ is divisible
 by $\varpi$, contradicting minimality.

 This saturation argument removes the entire initial vertical
 denominator, regardless of its dependence on $N_1$ and the degrees.
 Evaluation at the integral ordinary CM object now bounds
 $F_P(\mathrm{CM})=L_\phi(P)^2$.  Taking a scalar square root proves
 $|L_\phi(P)|_p\leq C$ and hence the required polynomial-norm estimate.

 \paragraph{The measure and its unit restriction.}
 We have obtained a bounded linear functional on polynomials on the
 compact integral position lattice.  Polynomial density extends it to
 continuous functions, hence to a bounded measure.  Approximate the
 characteristic function of the units by polynomials.  Applying the
 same square and cross-product argument to the difference proves that
 cutting this measure corresponds exactly to cutting the finite test.
 On the complement of the product unit locus, the monomials tend
 uniformly to zero when every coordinate degree tends to infinity.
 On units, after fixing the torsion component, the formula
 $x^s=\omega(x)^s\exp(s\log\langle x\rangle)$ proves analyticity on a
 sufficiently narrow branch.  The finitely many fixed coefficient-field
 unit scales have the same property after narrowing the branch once.
\end{proof}

\subsection{Mixed Fourier--Jacobi expansions and their integrality}

At the rank-one boundary use $x\in A$ for the hyperbolic variables and
the Schr\"odinger realization of $A\otimes B$ for the remaining variables.
The Fourier--Jacobi indices are the positive norms $(x,x)$, together
with zero.  At zero Jacobi argument a power of the positive hyperbolic
Fock column inserts the corresponding power of the coordinates of $x$.
There are no lower Hermite terms: the used creations are the isotropic
half of the variables for the contractions, namely holomorphic linear
coordinates on the positive trace form of $A_\infty$.
The positive and conjugate multiplicative frames in this direction are
toric logarithmic frames, so these insertions incur no CM periods.
The $B$-columns have precisely the unary normalization of
Lemma~\ref{lem:unary-moments}.

At a rational Jacobi torsion point, use Heisenberg translation of the
finite tests and its translated Fourier trivialization.  A rational
unipotent moves this point to zero; for prime-to-$p$ torsions this is a
prime-to-$p$ isogeny move and supplies a splitting of the semiabelian
differentials.  Consequently the same assertion holds monomial by
monomial according to column degree.  No estimate in an unrelated
analytic trivialization at a nonzero Jacobi argument is required.
Restricting to $B\oplus\mathbb H$ gives another description: after
removing $\exp(2\pi i(x,x)\tau)$, the hyperbolic coefficient is the
harmonic position polynomial.  In particular, at index zero only the
minor using the two $B$-columns remains; it is the binary theta lift of
parallel determinant weight $m$.

\begin{lemma}\label{lem:mixed-expansions}
 The mixed Fourier--Jacobi expansions of
 \eqref{eq:upper-theta-integral}, in algebraic ordinary differential
 frames and in the monomial-dual lattice, have a common bounded
 denominator.  The bound is independent of $m$, $i$, and the tame
 character, and is coefficientwise in the tame group ring.  The same
 assertion holds after substitution into any fixed $R_b$.
\end{lemma}
\begin{proof}
 Choose source integration representatives integral at $p$, by weak
 approximation in the rational unitary group, or by changing the
 rational Hermitian basis.  The source quotient then gives finite sums
 at a level of fixed volume denominator.  This remains true at $q_i$
 because the determinant equivariance of
 \eqref{eq:q-pivot-test} cancels the varying source character.
 The archimedean transformation factors cancel in the integral, and
 character evaluations on the remaining representatives are algebraic
 $p$-units.

 A boundary representative may be chosen integral at $p$ by Iwasawa
 decomposition, sending the specified multiplicative member to
 $e+B$ modulo the fixed level.  The translated $B$-tests remain
 integrally supported.  The test in the column Fourier-dualized to
 produce $x$ is invariant under integral translations, so $x$ too is
 integral at $p$.  Rational Jacobi torsions can use unipotent
 parameters divisible by the fixed $p$-level; these torsions are still
 dense and preserve this structure.  Representatives and unipotents
 can simultaneously be chosen integral at the moving and fixed
 exceptional places, using approximation in the cusp parabolic and
 the rational Cayley chart on its binary unitary factor.

 After the source-equivariance cancellation just made, expand each
 remaining finite test in the tame group basis and decompose each
 scalar coefficient test into simple unary tensors, retaining integral
 support and the original coefficient bound.  Apply the scalar unary
 estimate before reassembling the group-algebra coefficients.  Additive
 Fourier transforms away from $p$ cost only $p$-units; at $p$ the
 levels are fixed.  Each index coefficient at a prime-to-$p$ Jacobi
 torsion is consequently a finite sum of the bounded moments in
 Lemma~\ref{lem:unary-moments}, multiplied by integral powers of the
 coordinates of $x$ and by the unit scales already separated in
 Convention~\ref{conv:scaled-support}.  The nonarchimedean bound is
 unaffected by the number of summands.

 At each fixed weight and Fourier--Jacobi index, these algebraic
 evaluations first give an algebraic coefficient section.  Indeed the
 coefficient bundle on the abelian chart is algebraic.  By proper
 algebraic--analytic comparison, the holomorphic coefficient lies in
 the complexification of its finite-dimensional algebraic space of
 sections.  Choose finitely many algebraic
 prime-to-$p$ torsion evaluations that separate this space, in the
 algebraic split differential frames just used.  In an algebraic
 basis their evaluation matrix has full rank.  Algebraicity of the
 values therefore gives algebraicity of the section, with an initial
 denominator that may depend on the weight, index, and tame level.

 The full family of these torsion evaluations detects vertical integrality.
 The prime-to-$p$ torsions, and those
 satisfying the fixed congruences just imposed, are dense on the
 special fiber.  Their split differential frames are obtained by
 prime-to-$p$ isogenies.  After clearing an individual denominator, a
 section whose normalized torsion evaluations vanish modulo the
 uniformizer is therefore zero modulo it.  The same least-denominator
 argument as above proves the uniform bound for the coefficient
 section.  These algebraic coefficients first give algebraicity of the
 form by the characteristic-zero expansion principle.  The ordinary
 Fourier--Jacobi expansion principle then transfers their uniform
 bound to the form itself.
 The canonical cusp charts with their finite \'etale level are those
 of Lemma~\ref{lem:canonical-ordinary-level}.  Their construction uses
 only ordinary PEL geometry and the flag and frame data of $J$, so it
 is independent of the theta estimates and the cusp-lifting conclusion.

 Finally the substitution defining $R_b$ uses only the finitely many
 binomial coefficients through degree $b-1$.  Its factorial loss is a
 fixed constant depending on $b$.  This proves the last assertion.
\end{proof}

\begin{proof}[Proof of Proposition~\ref{prop:theta-family}]
 The Fock calculation gives the weight and holomorphy, the mixed
 expansion gives its extension at cusps and its algebraic integral
 interpretation, Lemma~\ref{lem:split-u-operators} gives the exact
 operator values, and Lemma~\ref{lem:theta-eigenparameters} gives the
 spherical parameters.  Lemma~\ref{lem:mixed-expansions} gives the
 uniform bounds.
\end{proof}
